\chapter{Rank sums and the canonical scan}

\begin{definition}\leanok
[Rank sum]
  \label{def:rank_sum}
  \lean{Ninegates.rankSum}
  \uses{def:counts}
  % SOURCE: half-formal.qmd, “点数和的模 3 不变量” (read from references/half-formal.qmd)
  % SOURCE QUOTE: “对长度为 $m$ 的重数向量 $d$，定义它的**点数和** $\operatorname{rankSum}(d)=\sum_{i=1}^m i\,d_i$。”
  \textit{Source: half-formal.qmd, rank-sum invariant section.}
  For a count vector \(d\) of length \(m\), define
  \[
    \operatorname{rankSum}(d)=\sum_{i=1}^{m}i\,d_i.
  \]
\end{definition}

\begin{definition}\leanok
[Scan state]
  \label{def:scan_state}
  \lean{Ninegates.ScanState}
  % SOURCE: half-formal.qmd, #def-canonical-scan (read from references/half-formal.qmd)
  % SOURCE QUOTE: “扫描状态由两个数 $(a,b)\in\{0,1,2\}^2$ 组成，其中： - $a$ 是从前两点开始、将在当前点使用一张牌的顺子数； - $b$ 是从前一点开始、将在当前点使用一张牌的顺子数。”
  \textit{Source: half-formal.qmd, Definition \texttt{\#def-canonical-scan}.}
  A scan state is a pair \((a,b)\in\{0,1,2\}^2\).  The entries count
  runs started two and one positions earlier that must consume the current
  rank.
\end{definition}

\begin{definition}\leanok
[One scan step]
  \label{def:scan_step}
  \lean{Ninegates.scanStep}
  \uses{def:scan_state}
  % SOURCE: half-formal.qmd, #def-canonical-scan (read from references/half-formal.qmd)
  % SOURCE QUOTE: “在当前坐标读到重数 $z$ 时： 1. 若 $a+b>z$，扫描失败； 2. 若 $a+b\le z$，令 $$u=z-a-b,\qquad s=u\bmod3,$$ 其中 $s\in\{0,1,2\}$ 是最小非负余数，并把状态更新为 $$(a,b)\longmapsto(b,s).$$”
  \textit{Source: half-formal.qmd, Definition \texttt{\#def-canonical-scan}.}
  On input multiplicity \(z\), the step from \((a,b)\) fails when
  \(a+b>z\).  Otherwise put \(u=z-a-b\), let \(s\) be the least
  nonnegative residue of \(u\) modulo \(3\), and return \((b,s)\).
\end{definition}

\begin{definition}\leanok
[Scanning a count vector]
  \label{def:scan}
  \lean{Ninegates.scan}
  \uses{def:counts,def:scan_state,def:scan_step}
  % SOURCE: half-formal.qmd, #def-canonical-scan (read from references/half-formal.qmd)
  % SOURCE QUOTE: “扫描从状态 $(0,0)$ 开始。读完全部坐标后，若状态回到 $(0,0)$，称扫描接受 $d$。”
  \textit{Source: half-formal.qmd, Definition \texttt{\#def-canonical-scan}.}
  Scan a count vector from left to right by repeated application of the
  one-step transition.  The result records either the first failure or the
  final state.  The operation may also be started from an arbitrary state,
  which is needed when a middle block is removed.
\end{definition}

\begin{definition}\leanok
[Scan acceptance]
  \label{def:scan_accepts}
  \lean{Ninegates.ScanAccepts}
  \uses{def:scan,def:scan_state}
  A count vector is accepted when its scan starts at \((0,0)\), never
  fails, and finishes at \((0,0)\).
\end{definition}

\begin{lemma}

\leanok
[Singleton finite sum by finite value]
  \label{lem:sum_if_fin_val_eq}
  \lean{Ninegates.sum_if_fin_val_eq}
  For a function on a finite interval, the sum that keeps only entries
  whose index has natural value \(i\) is the corresponding entry when
  \(i\) is in range, and is \(0\) otherwise.
\end{lemma}
\begin{proof}

\leanok
  If \(i\) is in range, exactly one finite index has that natural value,
  so the sum has one nonzero term.  If \(i\) is out of range, no index can
  satisfy the equality, so every summand is zero.
\end{proof}

\begin{lemma}

\leanok
[Weight of one run in the rank sum]
  \label{lem:run_weight_sum}
  \lean{Ninegates.run_weight_sum}
  \uses{def:rank_sum,lem:sum_if_fin_val_eq}
  A run beginning at zero-based start \(j\) contributes the three
  one-based ranks \((j+1),(j+2),(j+3)\).  Hence its weighted contribution
  is \(3(j+2)\) times the run multiplicity.
\end{lemma}
\begin{proof}

\leanok
  Split the finite sum into the three possible coordinates of the run.
  The singleton-sum lemma evaluates those three terms, and their weights
  add to \((j+1)+(j+2)+(j+3)=3(j+2)\).
\end{proof}

\begin{lemma}

\leanok
[Finite sum of zero congruence terms]
  \label{lem:sum_modEq_zero}
  \lean{Ninegates.sum_modEq_zero}
  If every term in a finite sum is congruent to \(0\) modulo \(3\), then
  the whole sum is congruent to \(0\) modulo \(3\).
\end{lemma}
\begin{proof}

\leanok
  Induct over the finite set of summation indices.  The empty sum is zero,
  and adjoining one term preserves congruence to zero under addition.
\end{proof}

\begin{lemma}\leanok
[Rank sum of a meldable vector]
  \label{lem:meldable_rank_sum}
  \lean{Ninegates.meldable_rankSum_mod}
  \uses{def:meldable,def:rank_sum,lem:run_weight_sum,lem:sum_modEq_zero}
  % SOURCE: half-formal.qmd, #lem-rank-sum (read from references/half-formal.qmd)
  % SOURCE QUOTE: “若 $d$ 可面子分解，则 $\operatorname{rankSum}(d)\equiv0\pmod3$。”
  \textit{Source: half-formal.qmd, Lemma \texttt{\#lem-rank-sum}.}
  If \(d\) is meldable, then
  \(\operatorname{rankSum}(d)\equiv0\pmod3\).
\end{lemma}
% SOURCE QUOTE PROOF: “设 $d$ 的一个面子分解由刻子数 $t_i$ 和顺子数 $r_i$ 给出。一个点数 $i$ 的刻子对点数和的贡献是 $3i$，它能被 $3$ 整除。一个从点数 $i$ 开始的顺子贡献 $i+(i+1)+(i+2)=3i+3$，也能被 $3$ 整除。把所有刻子与顺子的贡献相加，得到 $$ \operatorname{rankSum}(d) =\sum_i 3i\,t_i+\sum_i(3i+3)r_i \equiv0\pmod3.$$”
\begin{proof}

\leanok
  Choose triplet counts \(t_i\) and run counts \(r_i\) witnessing
  meldability.  A triplet at \(i\) contributes \(3i\) to the rank sum,
  while a run starting at \(i\) contributes
  \(i+(i+1)+(i+2)=3i+3\).  Every contribution is divisible by \(3\);
  summing them proves the congruence.
\end{proof}

\begin{lemma}

\leanok
[Rank sum is additive]
  \label{lem:rankSum_add}
  \lean{Ninegates.rankSum_add}
  \uses{def:rank_sum}
  The rank sum of a pointwise sum of two count vectors equals the sum of
  their rank sums.
\end{lemma}
\begin{proof}

\leanok
  Expand the weighted finite sum and distribute multiplication over
  coordinatewise addition, then separate the finite sum into the two
  component sums.
\end{proof}

\begin{lemma}

\leanok
[Rank sum of two basis tiles]
  \label{lem:rankSum_two_basis}
  \lean{Ninegates.rankSum_two_basis}
  \uses{def:rank_sum,def:basis}
  Two copies of the basis vector at \(p\) have rank sum equal to twice the
  one-based rank of \(p\).
\end{lemma}
\begin{proof}

\leanok
  All coordinates except \(p\) vanish.  The remaining coordinate has
  multiplicity \(2\), so the weighted sum is \(2(p+1)\).
\end{proof}

\begin{lemma}

\leanok
[Rank sum of one basis tile]
  \label{lem:rankSum_basis}
  \lean{Ninegates.rankSum_basis}
  \uses{def:rank_sum,def:basis}
  The basis vector at \(p\) has rank sum equal to the one-based rank of
  \(p\).
\end{lemma}
\begin{proof}

\leanok
  The basis vector has exactly one nonzero coordinate, at \(p\), with
  multiplicity \(1\).  The weighted sum is therefore \(p+1\).
\end{proof}

\begin{lemma}

\leanok
[Adding one basis tile]
  \label{lem:rankSum_add_basis}
  \lean{Ninegates.rankSum_add_basis}
  \uses{lem:rankSum_add,lem:rankSum_basis}
  Adding one tile at \(x\) increases the rank sum by the one-based rank
  of \(x\).
\end{lemma}
\begin{proof}

\leanok
  Apply additivity of the rank sum to \(c+e_x\), then evaluate the rank
  sum of the basis vector.
\end{proof}

\begin{lemma}

\leanok
[Rank sum after removing a pair]
  \label{lem:removePair_rankSum_add_pair}
  \lean{Ninegates.removePair_rankSum_add_pair}
  \uses{def:remove_pair,def:basis,lem:rankSum_add,lem:rankSum_two_basis}
  If \(w_p\geq2\), then the rank sum of \(w-2e_p\), plus twice the
  one-based rank of \(p\), is the rank sum of \(w\).
\end{lemma}
\begin{proof}

\leanok
  Coordinatewise, adding back two basis tiles at \(p\) to the pair-removed
  vector recovers \(w\).  Additivity and the two-basis evaluation then
  give the displayed rank-sum identity.
\end{proof}

\begin{lemma}\leanok
[Rank sum after choosing the pair]
  \label{lem:complete_rank_sum}
  \lean{Ninegates.complete_rankSum_mod}
  \uses{def:basis,def:remove_pair,def:meldable,def:rank_sum,lem:meldable_rank_sum,lem:rankSum_add_basis,lem:removePair_rankSum_add_pair}
  % SOURCE: half-formal.qmd, #lem-rank-sum (read from references/half-formal.qmd)
  % SOURCE QUOTE: “进一步，设 $c$ 是长度为 $n$ 的牌型，并记 $S=\operatorname{rankSum}(c)$。如果进张 $x$ 后，$c+e_x$ 存在一个以点数 $p$ 为雀头的和牌分解，则 $S+x\equiv2p\pmod3$。”
  \textit{Source: half-formal.qmd, Lemma \texttt{\#lem-rank-sum}.}
  If \(c+e_x\) has a complete decomposition whose pair is at \(p\),
  then
  \[
    \operatorname{rankSum}(c)+x\equiv2p\pmod3.
  \]
\end{lemma}
% SOURCE QUOTE PROOF: “现在设 $c+e_x$ 以 $p$ 为雀头。按照 @def-complete， $d=c+e_x-2e_p$ 可面子分解，故 $\operatorname{rankSum}(d)\equiv0\pmod3$。另一方面，点数和对向量加减是线性的，所以 $\operatorname{rankSum}(d)=S+x-2p$。因此 $S+x-2p\equiv0\pmod3$，等价地，$S+x\equiv2p\pmod3$。”
\begin{proof}

\leanok
  Removing the chosen pair gives the meldable vector
  \(d=c+e_x-2e_p\).  Its rank sum is divisible by \(3\) by
  Lemma~\ref{lem:meldable_rank_sum}.  Directly expanding the finite sum
  gives
  \(\operatorname{rankSum}(d)=\operatorname{rankSum}(c)+x-2p\);
  rearranging modulo \(3\) gives the claim.
\end{proof}

\begin{definition}

\leanok
[Extended start lookup]
  \label{def:scan_start_at}
  \lean{Ninegates.scanStartAt}
  \uses{def:counts}
  Given normalized run-start counts, extend them to all natural indices by
  returning the indexed count in range and \(0\) out of range.
\end{definition}

\begin{definition}

\leanok
[Cons for count vectors]
  \label{def:counts_cons}
  \lean{Ninegates.countsCons}
  \uses{def:counts}
  A count vector can be obtained by placing a head value before a tail
  vector, with indices described by their natural values.
\end{definition}

\begin{definition}

\leanok
[First incoming run component]
  \label{def:scan_prev2}
  \lean{Ninegates.scanPrev2}
  \uses{def:scan_state,def:counts,def:scan_start_at}
  At coordinate \(i\), the first incoming scan component is the initial
  state's first component at \(i=0\), the initial state's second component
  at \(i=1\), and otherwise the start count at \(i-2\).
\end{definition}

\begin{definition}

\leanok
[Second incoming run component]
  \label{def:scan_prev1}
  \lean{Ninegates.scanPrev1}
  \uses{def:scan_state,def:counts,def:scan_start_at}
  At coordinate \(i\), the second incoming scan component is the initial
  state's second component at \(i=0\), and otherwise the start count at
  \(i-1\).
\end{definition}

\begin{definition}

\leanok
[First final scan component]
  \label{def:scan_out1}
  \lean{Ninegates.scanOut1}
  \uses{def:scan_state,def:counts}
  The first component of the final scan state is the initial first
  component for the empty vector, the initial second component for length
  one, and the penultimate start count for longer vectors.
\end{definition}

\begin{definition}

\leanok
[Second final scan component]
  \label{def:scan_out2}
  \lean{Ninegates.scanOut2}
  \uses{def:scan_state,def:counts}
  The second component of the final scan state is the initial second
  component for the empty vector, and the final start count for nonempty
  vectors.
\end{definition}

\begin{lemma}

\leanok
[Output shift after adding a head, first component]
  \label{lem:scan_out1_cons}
  \lean{Ninegates.scanOut1_cons}
  \uses{def:scan_out1,def:counts_cons}
  If a first scan step changes the initial state from \(\sigma\) to
  \(\sigma'\), then adding that head start before a tail preserves the
  computed final first component.
\end{lemma}
\begin{proof}

\leanok
  Check the possible lengths of the tail.  For length zero and one the
  claim is exactly the relation between \(\sigma\) and \(\sigma'\); for
  longer tails both sides select the same shifted start count.
\end{proof}

\begin{lemma}

\leanok
[Output shift after adding a head, second component]
  \label{lem:scan_out2_cons}
  \lean{Ninegates.scanOut2_cons}
  \uses{def:scan_out2,def:counts_cons}
  Adding a head start before a tail preserves the computed final second
  component after updating the scan state.
\end{lemma}
\begin{proof}

\leanok
  By cases on the tail length, the formula either reads the updated
  initial component or the same final shifted start count.
\end{proof}

\begin{lemma}

\leanok
[Output shift after removing a head, first component]
  \label{lem:scan_out1_tail}
  \lean{Ninegates.scanOut1_tail}
  \uses{def:scan_out1}
  Removing the head of the normalized starts shifts the first final output
  component exactly as a successful first scan step does.
\end{lemma}
\begin{proof}

\leanok
  Compare the formulas by the remaining length.  The boundary cases read
  the updated initial state, and the longer case reads the same shifted
  penultimate start count.
\end{proof}

\begin{lemma}

\leanok
[Output shift after removing a head, second component]
  \label{lem:scan_out2_tail}
  \lean{Ninegates.scanOut2_tail}
  \uses{def:scan_out2}
  Removing the head of the normalized starts shifts the second final
  output component exactly as a successful first scan step does.
\end{lemma}
\begin{proof}

\leanok
  The empty tail reads the updated second component.  A nonempty tail reads
  the same shifted final start count on both sides.
\end{proof}

\begin{lemma}

\leanok
[Incoming first component under cons]
  \label{lem:scan_prev2_cons_succ}
  \lean{Ninegates.scanPrev2_cons_succ}
  \uses{def:scan_prev2,def:scan_start_at,def:counts_cons}
  After adding a head start, the first incoming component at a successor
  coordinate agrees with the corresponding first incoming component in the
  tail scan.
\end{lemma}
\begin{proof}

\leanok
  Split the successor coordinate into the first two boundary positions and
  the remaining positions.  In the remaining case the extended lookup
  shifts from \(i-2\) in the longer vector to the corresponding tail
  index.
\end{proof}

\begin{lemma}

\leanok
[Incoming second component under cons]
  \label{lem:scan_prev1_cons_succ}
  \lean{Ninegates.scanPrev1_cons_succ}
  \uses{def:scan_prev1,def:scan_start_at,def:counts_cons}
  After adding a head start, the second incoming component at a successor
  coordinate agrees with the corresponding second incoming component in
  the tail scan.
\end{lemma}
\begin{proof}

\leanok
  The first successor coordinate is the boundary case reading the updated
  state.  All later coordinates read the same shifted start count.
\end{proof}

\begin{lemma}

\leanok
[Incoming first component under tail]
  \label{lem:scan_prev2_tail}
  \lean{Ninegates.scanPrev2_tail}
  \uses{def:scan_prev2,def:scan_start_at}
  Removing a head start shifts the first incoming component on every tail
  coordinate to the first incoming component at the successor coordinate
  of the original scan.
\end{lemma}
\begin{proof}

\leanok
  The first two tail coordinates are the boundary components of the
  updated state.  Every later coordinate reduces to the same shifted
  extended start lookup.
\end{proof}

\begin{lemma}

\leanok
[Incoming second component under tail]
  \label{lem:scan_prev1_tail}
  \lean{Ninegates.scanPrev1_tail}
  \uses{def:scan_prev1,def:scan_start_at}
  Removing a head start shifts the second incoming component on every tail
  coordinate to the second incoming component at the successor coordinate
  of the original scan.
\end{lemma}
\begin{proof}

\leanok
  The first coordinate reads the updated state component; every later
  coordinate reads the same shifted extended start lookup.
\end{proof}

\begin{definition}

\leanok
[General scan decomposition]
  \label{def:scan_decomp}
  \lean{Ninegates.ScanDecomp}
  \uses{def:counts,def:scan_state,def:scan_prev2,def:scan_prev1,def:scan_out1,def:scan_out2}
  A decomposition from initial state \(\sigma\) to final state \(\tau\)
  consists of triplet counts and normalized start counts \(s_i<3\), such
  that \(\tau\) is the state predicted by the starts and every coordinate
  equals the two incoming starts, the new start, and three times a triplet
  count.
\end{definition}

\begin{lemma}

\leanok
[One scan step from a coordinate decomposition]
  \label{lem:scan_step_of_coord}
  \lean{Ninegates.scan_step_of_coord}
  \uses{def:scan_step}
  If a coordinate has the form \(a+b+s+3t\), where \((a,b)\) is the
  incoming state and \(s<3\), then one scan step succeeds and outputs
  \((b,s)\).
\end{lemma}
\begin{proof}

\leanok
  The coordinate is at least \(a+b\), so the step cannot fail.  After
  subtracting the incoming components, the remainder is \(s+3t\), whose
  residue modulo \(3\) is \(s\).
\end{proof}

\begin{lemma}

\leanok
[Coordinate decomposition from one scan step]
  \label{lem:coord_of_scan_step}
  \lean{Ninegates.coord_of_scan_step}
  \uses{def:scan_step}
  A successful scan step from \(\sigma\) on input \(z\) to \(\tau\)
  provides a triplet count \(t\) such that
  \(z=\sigma_1+\sigma_2+\tau_2+3t\), and the first component of
  \(\tau\) is the second component of \(\sigma\).
\end{lemma}
\begin{proof}

\leanok
  Success means \(z\) is at least the two incoming components.  Divide the
  residual by \(3\); the quotient is \(t\), and the remainder is exactly
  the new run-start component.
\end{proof}

\begin{lemma}

\leanok
[Scan equals generalized decomposition]
  \label{lem:scan_iff_scan_decomp}
  \lean{Ninegates.scan_iff_scanDecomp}
  \uses{def:scan,def:scan_decomp,def:counts_cons,lem:scan_step_of_coord,lem:coord_of_scan_step,lem:scan_out1_cons,lem:scan_out2_cons,lem:scan_out1_tail,lem:scan_out2_tail,lem:scan_prev2_cons_succ,lem:scan_prev1_cons_succ,lem:scan_prev2_tail,lem:scan_prev1_tail}
  A scan of a vector from \(\sigma\) to \(\tau\) succeeds if and only if
  the vector admits the generalized coordinate decomposition from
  \(\sigma\) to \(\tau\).
\end{lemma}
\begin{proof}

\leanok
  Induct on the vector length.  The empty case records that the initial
  and final states coincide.  For a nonempty vector, a successful first
  scan step gives the head triplet count and head start; the induction
  hypothesis handles the tail, and the shift lemmas translate tail
  coordinates and final outputs back to the original vector.  Conversely,
  the head coordinate decomposition forces the first scan step to succeed,
  and the tail shift lemmas provide the decomposition needed by the
  induction hypothesis.
\end{proof}

\begin{lemma}

\leanok
[Run cover at one coordinate]
  \label{lem:run_cover_at}
  \lean{Ninegates.run_cover_at}
  \uses{lem:sum_if_fin_val_eq}
  For any coordinate \(i\), the finite sum over all run starts whose
  length-three window covers \(i\) has only the three possible
  contributors starting at \(i\), \(i-1\), and \(i-2\), with any
  out-of-range contributor interpreted as zero.
\end{lemma}
\begin{proof}

\leanok
  A start \(j\) covers \(i\) precisely when \(j\leq i\leq j+2\), which is
  equivalent to \(j\in\{i,i-1,i-2\}\).  Split the sum into these singleton
  possibilities and use the out-of-range case of the singleton-sum lemma
  for nonexistent starts.
\end{proof}

\begin{theorem}\leanok
[Canonical scan characterizes meldability]
  \label{thm:meldable_iff_scan}
  \lean{Ninegates.meldable_iff_scanAccepts}
  \uses{def:meldable,def:scan_accepts,def:scan_decomp,lem:scan_iff_scan_decomp,lem:run_cover_at}
  % SOURCE: half-formal.qmd, #lem-recurrence (read from references/half-formal.qmd)
  % SOURCE QUOTE: “非负整数向量 $d=(d_1,\ldots,d_m)$ 可面子分解，当且仅当 @def-canonical-scan 中的规范扫描接受 $d$。”
  \textit{Source: half-formal.qmd, Lemma \texttt{\#lem-recurrence}.}
  A count vector is meldable if and only if its canonical scan is
  accepted.
\end{theorem}
% SOURCE QUOTE PROOF: “假设 $d$ 可面子分解。于是存在刻子计数 $t_i$ 与顺子计数 $r_i$，满足 @def-meldable 中的逐坐标等式。首先把这个分解规范化。若某个 $r_i\ge3$，则三组相同的顺子 $(i,i+1,i+2)$ 可以替换为三个刻子 $(i,i,i)$、$(i+1,i+1,i+1)$、$(i+2,i+2,i+2)$。这个替换不改变任何点数的总张数，并把 $r_i$ 减少 $3$。对每个顺子起点 $i$ 重复这一操作，最终得到另一个分解，使得 $0\le r_i\le2$ 对所有 $i$ 成立。这个过程必然终止，因为每次替换都会使非负整数 $\sum_i r_i$ 减少 $3$。现在从左到右考察这个规范分解。在点数 $i$ 处，从 $i-2$ 与 $i-1$ 开始的顺子已经确定占用 $r_{i-2}+r_{i-1}$ 张点数 $i$。因为这是一个合法分解，必有 $r_{i-2}+r_{i-1}\le d_i$。去掉这些已经承诺给早先顺子的牌，剩余牌数为 $u_i=d_i-r_{i-2}-r_{i-1}$。剩余牌只能用于两类面子：点数 $i$ 的刻子，以及从 $i$ 开始的顺子。因此 $u_i=3t_i+r_i$。又因为 $0\le r_i\le2$，所以 $r_i$ 必定等于 $u_i$ 除以 $3$ 的最小非负余数。也就是说，只要扫描在此前各坐标产生了 $s_{i-2}=r_{i-2}$ 与 $s_{i-1}=r_{i-1}$，它在当前坐标就必然产生 $s_i=r_i$。从初始条件 $s_{-1}=s_0=r_{-1}=r_0=0$ 出发，用归纳法得到扫描在每一步都不会失败，而且对所有 $i$ 都有 $s_i=r_i$。顺子不能从点数 $m-1$ 或 $m$ 开始，因此按下标约定， $r_{m-1}=r_m=0$。于是 $s_{m-1}=s_m=0$，扫描接受 $d$。反过来，假设扫描接受 $d$。对每个 $i$，扫描未失败保证 $u_i=d_i-s_{i-2}-s_{i-1}\ge0$。由于 $s_i$ 是 $u_i$ 除以 $3$ 的余数， $u_i-s_i$ 是非负的 $3$ 的倍数。因此 $t_i=(u_i-s_i)/3$ 是非负整数。整理 $u_i$ 的定义可得 $$d_i=3t_i+s_i+s_{i-1}+s_{i-2}. \tag{2}$$ 扫描接受还给出 $s_{m-1}=s_m=0$。因此可以把 $s_i$ 视为从 $i$ 开始的顺子数：在最后两个点数处没有顺子越出范围。式 (2) 正是 @def-meldable 所要求的逐坐标等式，所以 $t_i$ 与 $s_i$ 构成 $d$ 的一个合法面子分解。”
\begin{proof}

\leanok
  First suppose a meld decomposition is given.  At each run start, replace
  groups of three identical runs by one triplet at each of their three
  ranks.  This preserves the count vector and leaves every run count in
  \(\{0,1,2\}\).  Inducting from left to right, the two incoming run counts
  are exactly the scan state.  After subtracting them, the residual is
  \(3t_i+r_i\), so its residue modulo \(3\) is the normalized run count
  \(r_i\).  The scan therefore follows these run counts, never fails, and
  ends at zero because no run starts in either of the final two positions.

  Conversely, suppose the scan accepts.  At coordinate \(i\), let \(s_i\)
  be the new second state component and let
  \(u_i=d_i-s_{i-2}-s_{i-1}\).  Nonfailure gives \(u_i\geq0\), and the
  transition rule makes \(u_i-s_i\) a nonnegative multiple of \(3\).
  Set \(t_i=(u_i-s_i)/3\).  Then
  \[
    d_i=3t_i+s_i+s_{i-1}+s_{i-2}.
  \]
  Acceptance says the final two run counts are zero, so these \(t_i,s_i\)
  form a legal meld decomposition.

  In the finite-vector implementation, package the recurrence as the
  generalized decomposition from Lemma~\ref{lem:scan_iff_scan_decomp}.  In
  the forward direction, replace every original run count \(r_j\) by its
  residue \(r_j\bmod 3\) plus \(3\lfloor r_j/3\rfloor\); the quotient
  contributions are absorbed into the triplet counts.  In the reverse
  direction, Lemma~\ref{lem:run_cover_at} rewrites the original meldability
  cover sum at coordinate \(i\) as exactly the incoming starts from
  \(i-2\), \(i-1\), and \(i\).  The final two output equalities ensure
  that no run starts in the two invalid right-end positions.
\end{proof}

\begin{lemma}\leanok
[Reading one singleton]
  \label{lem:scan_step_one}
  \lean{Ninegates.scanStep_one}
  \uses{def:scan_state,def:scan_step}
  % SOURCE: half-formal.qmd, #lem-three-ones-state (read from references/half-formal.qmd)
  % SOURCE QUOTE: “规范扫描读入一个重数为 $1$ 的坐标时，可行输入状态恰好是 $(0,0),(0,1),(1,0)$，相应的状态转移为 $$(0,0)\longmapsto(0,1),\qquad (0,1)\longmapsto(1,0),\qquad (1,0)\longmapsto(0,0).$$”
  \textit{Source: half-formal.qmd, Lemma \texttt{\#lem-three-ones-state}.}
  A scan step on input \(1\) succeeds exactly from
  \((0,0),(0,1),(1,0)\), and maps these states respectively to
  \((0,1),(1,0),(0,0)\).
\end{lemma}
% SOURCE QUOTE PROOF: “设当前输入状态为 $(a,b)$，当前重数为 $1$。可行性要求 $a+b\le1$。又因为规范扫描中的状态分量属于 $\{0,1,2\}$，满足此不等式的状态只有 $(0,0),(0,1),(1,0)$。逐一代入 @def-canonical-scan： - 从 $(0,0)$ 出发，$u=1$，新顺子数为 $1\bmod3=1$，故输出 $(0,1)$； - 从 $(0,1)$ 出发，$u=0$，新顺子数为 $0$，故输出 $(1,0)$； - 从 $(1,0)$ 出发，$u=0$，新顺子数为 $0$，故输出 $(0,0)$。”
\begin{proof}

\leanok
  Success requires \(a+b\leq1\).  Since both components lie in
  \(\{0,1,2\}\), this leaves exactly the three displayed states.
  Substitution in the transition formula gives the stated three-cycle.
\end{proof}

\begin{lemma}\leanok
[Three singletons preserve scan state]
  \label{lem:scan_three_ones}
  \lean{Ninegates.scan_three_ones}
  \uses{def:scan,lem:scan_step_one}
  % SOURCE: half-formal.qmd, #lem-three-ones-state (read from references/half-formal.qmd)
  % SOURCE QUOTE: “因此，如果从某一状态开始能够连续读入三个重数 $1$，那么读完后回到原状态。换言之，单词 $(1,1,1)$ 在其可行状态域上作用为恒等映射。”
  \textit{Source: half-formal.qmd, Lemma \texttt{\#lem-three-ones-state}.}
  If a scan starting from state \(\sigma\) can read three consecutive
  inputs \(1,1,1\), its state after those inputs is again \(\sigma\).
\end{lemma}
% SOURCE QUOTE PROOF: “这三个状态组成一个长度为 $3$ 的循环，所以连续应用三次后回到原状态。”
\begin{proof}

\leanok
  By Lemma~\ref{lem:scan_step_one}, every successful state belongs to the
  displayed three-cycle.  Applying that cycle three times is the identity.
\end{proof}
